%% kapa: exam=1 round=27 year=2016 period=1 form=A subject=civil_law subject_ko=민법 seq=16 number=16
%% source: Q-Net 감정평가사 기출문제 (cst00309) artlSeq=5169402 file=extracted/1008/5169402/제27회 감정평가사 1차 문제/제27회 감정평가사 1교시문제(A형).pdf
%% conversion: kordoc 4.24.51 -> tools/gen1.py (auto); review=proofread-sonnet; flags=none
\begin{question}{16}
소멸시효의 기산점에 관한 설명으로 옳지 않은 것은? (다툼이 있으면 판례에 따름)
\begin{choices}
  \item 채무의 이행기가 도래한 후에 채무자의 요청에 의하여 채권자가 채무자에게 기한을 유예한 경우, 유예한 이행기일로부터 다시 소멸시효가 진행한다.
  \item 불확정기한부 채권은 객관적으로 기한이 도래하면 그 때부터 소멸시효가 진행한다.
  \item 동시이행의 항변권이 붙어있는 채권은 그 항변권이 소멸된 이후부터 소멸시효가 진행한다.
  \item 매수인이 매매 목적물인 부동산을 인도 받아 점유하고 있는 이상 그 소유권이전등기청구권의 소멸시효는 진행되지 않는다.
  \item 매매로 인한 소유권이전채무의 이행불능으로 인한 손해배상채권의 소멸시효는 그 소유권이전채무가 이행불능으로 된 때부터 진행한다.
\end{choices}
\end{question}
\answer{3}
